\documentclass[11pt]{article}

\usepackage[T1]{fontenc}
\usepackage[margin=1in]{geometry}
\usepackage{amsmath,amssymb,amsthm}
\usepackage{mathtools}
\usepackage{booktabs}
\usepackage{enumitem}
\usepackage{xcolor}
\usepackage{hyperref}
\usepackage{titlesec}
\usepackage[expansion=false]{microtype}

\hypersetup{
  colorlinks=true,
  linkcolor=black,
  urlcolor=blue!60!black,
  citecolor=black
}

\titleformat{\section}{\large\bfseries}{\thesection}{0.6em}{}
\titleformat{\subsection}{\normalsize\bfseries}{\thesubsection}{0.6em}{}

\newcommand{\Qhat}{\widehat{Q}}
\newcommand{\Eb}{\mathbb{E}}
\newcommand{\Rb}{\mathbb{R}}
\newcommand{\Lam}{\Lambda}
\newcommand{\Ktest}{K_{\mathrm{test}}}
\newcommand{\Con}{C_{\mathrm{on}}}
\newcommand{\epsQ}{\varepsilon_Q}

\newtheorem*{remark}{Remark}

\title{\bfseries Methodology Review\\[0.3em]
\large \emph{Large Language Models for Sequential Decision-Making:\\
Improving In-Context Learning via Supervised Fine-Tuning}\\[0.4em]
\normalsize Zhang, Aghaei \& Saghafian (arXiv:2605.09009v1)}
\author{}
\date{}

\begin{document}
\maketitle
\vspace{-2.5em}

\section{Summary of the contribution}

The paper proposes a supervised fine-tuning (SFT) recipe for endowing a pretrained,
open-source LLM (Llama-2-7B) with in-context sequential decision-making ability. The
pipeline is: (i) sample tasks $\tau$ from a task distribution $\mathcal{T}$; (ii) compute
a (near-)optimal oracle policy per task via backward induction; (iii) roll out
demonstration trajectories under that oracle; (iv) serialize trajectories into strings and
fine-tune via QLoRA on a next-token loss; (v) at test time present a few oracle support
trajectories in-context on \emph{held-out} tasks and act greedily on the model's predicted
action token. Three environment classes are studied---MDP (full observability), POMDP
(partial observability), APOMDP (model ambiguity via an $\alpha$-MEU robust criterion).
The headline empirical finding is that SFT roughly halves the optimality gap relative to
ICL-only and far exceeds a random floor, with the largest gains in long-horizon, low-fidelity,
and ambiguous regimes. A theoretical section interprets a single linear self-attention (LSA)
layer as an implicit in-context ridge-like estimator of the optimal $Q$-function in linear
MDPs and lifts a per-query prediction bound to an end-to-end policy suboptimality bound.

The work is competently executed and the writing is clear. This review concentrates on
methodological soundness rather than novelty, and is organized around (1) the
experimental design, (2) the evaluation protocol and statistics, (3) the theory--experiment
linkage, (4) the baselines, and (5) reproducibility. I flag concerns by severity throughout.

\section{Experimental design}

\subsection{Task generation and the train/test split}

The MDP is a $10$-state ($E=9$) ``energy-management'' chain with three actions, two of which
($a=0,a=2$) are deliberately redundant charge actions. Tasks differ only in a single scalar
transition probability $p\in[0.5,1)$ sampled uniformly; the POMDP adds an emission fidelity
$q\in\{0.5,0.8,1.0\}$; the APOMDP adds a KL-ball ambiguity set of size $|\mathcal{M}_\tau|\in\{1,3,5\}$.

\paragraph{Concern (moderate): the task family is very low-dimensional.} Across draws, an
MDP ``task'' is fully determined by one number $p$ on a bounded interval, with all other
structure (state count, reward shape $r=s/E$, action semantics, cost $c=-0.02$) held fixed.
This is a thin notion of ``generalization across tasks.'' A model that simply learns the
single induced policy family parameterized by $p$ can score well without any genuine
in-context task identification. The duplicated-action device is described as making ``action
selection harder,'' but it inflates the action space without adding decision-relevant
structure; it tests label disambiguation, not planning. The paper would be substantially
stronger with at least one task axis that varies reward structure or topology, so that the
few-shot context must actually \emph{identify} the task rather than index a one-parameter
sweep. As written, it is hard to distinguish in-context task inference from amortized
memorization of a one-dimensional policy manifold.

\paragraph{Concern (moderate): the split is over task realizations, not over structure.}
Train and test tasks are drawn from the same $\mathcal{T}$ with the same generative form;
``held-out'' means a fresh $p$ (or $q$, or ambiguity seed), not a held-out region of task
space. This is a legitimate i.i.d.\ generalization test but a weak one, and it should not be
read as evidence of compositional or structural transfer. The OOD experiments
(\S F.2) partially address this---testing $p\in[0.3,0.5)$, off-grid $q$, off-$\alpha$---and
are the most informative robustness evidence in the paper. However, the authors themselves
note the MDP OOD region is plausibly \emph{easier} (smaller $p$ $\Rightarrow$ less frequent
state change), which means the near-zero OOD gap is weak evidence of robustness rather than
strong evidence; an OOD direction known to be \emph{harder} would be more probative.

\subsection{The oracle and the ``optimal'' label}

For APOMDPs the supervision target is the $\alpha$-MEU robust-backward-induction policy.

\paragraph{Concern (moderate): ``optimal'' is overloaded, and the APOMDP optimality gap is
self-referential.} The optimality gap (Eq.\ 11) normalizes by $\mathrm{OPT}_\tau$, the reward
``under the optimal policy.'' In the APOMDP case the ``optimal'' policy is itself the
$\alpha$-MEU oracle the model was trained to imitate, evaluated under the same nominal
model used to construct it. The gap therefore measures \emph{imitation fidelity to the
$\alpha$-MEU oracle}, not optimality against the true (possibly mis-specified) data-generating
process. This is defensible as a behavioral-cloning metric but the paper's framing
(``handle model ambiguity'') suggests a robustness claim it does not actually test: there is
no evaluation under a genuinely adversarial or held-out member of the ambiguity set, which
is the entire point of robust control. The $\alpha=1$ result (larger $|\mathcal{M}|$
$\Rightarrow$ larger gap) is consistent with the oracle being harder to imitate, not
necessarily with the LLM being more or less robust.

\paragraph{Strength.} Using a model-based oracle as the supervision signal is the right call
for a synthetic study and is honestly flagged as a limitation (\S G) for real domains where
no oracle exists. The motivation---offline, oracle-labeled data in healthcare---is coherent.

\subsection{Serialization and the action-token decoding rule}

Trajectories are serialized as \texttt{<O\_t> 3, <A\_t> 1, <R\_t> 0.33, \dots} and the model
acts by restricting outputs to action tokens and taking the argmax.

\paragraph{Concern (minor but consequential): tokenization of numeric fields is unexamined.}
Rewards such as $0.33$ and multi-digit observations are fed as raw decimal strings to a
BPE tokenizer that was never designed for numeric regularity. How $0.33$ vs.\ $0.44$ tokenize,
and whether the model can read the reward channel at all, materially affects what ``in-context
learning'' is even possible here. None of this is ablated. Since the central claim is that the
context carries task information, the encoding of that context is a first-order design choice,
not an implementation detail. A digit-level or fixed-precision encoding ablation would close
this gap.

\paragraph{Concern (minor): greedy argmax over action tokens collapses the policy.} Restricting
to the single highest-probability action token discards the distributional output the loss
trained. For stochastic-optimal settings (and the POMDP belief-optimal action can be
genuinely mixed) this biases evaluation. The choice is reasonable for reproducibility but its
effect is unquantified.

\section{Evaluation protocol and statistics}

\subsection{Metric definition}

The optimality gap (Eq.\ 11) is a \emph{ratio} $\frac{\mathrm{OPT}_\tau - \mathrm{OPT}^{\mathrm{Eval}}_\tau}{\mathrm{OPT}_\tau}$ averaged over test tasks.

\paragraph{Concern (moderate): the normalization is fragile and possibly sign-ambiguous.}
Rewards here include a per-step cost $c=-0.02$ and a work reward $r=s/E\in[0,1]$. Cumulative
$\mathrm{OPT}_\tau$ can be small or, for charge-heavy optimal policies under some $p$, even
negative. Dividing by a near-zero or negative denominator makes the per-task ratio unstable
and the cross-task average ill-behaved; a single task with $\mathrm{OPT}_\tau\approx 0$ can
dominate. The paper reports clean percentages, so either denominators are comfortably bounded
away from zero in practice or outliers are being absorbed silently. This needs to be stated.
A more robust metric is a \emph{regret} (un-normalized reward difference) or a normalization
by $(\mathrm{OPT}_\tau - \mathrm{RANDOM}_\tau)$, which bounds the gap in $[0,1]$ and is
standard. The random ``floor'' is already computed, so this is cheap.

\paragraph{Concern (minor): the floor is reported inconsistently.} The MDP random policy gaps
(57.6\%, 49.6\%, 46.4\%) are well below the $100\%$ a normalized-to-random metric would give,
confirming the metric is gap-to-optimal-over-optimal, not gap-to-optimal-over-range. That is
fine, but it means ``SFT halves the gap'' is a statement about an unbounded-below quantity and
its magnitude is not directly comparable across $T$ or across settings with different reward
scales. Comparisons such as ``$15\%$ at $T=10$ vs.\ $15\%$ at $T=15$'' should not be read as
equal difficulty.

\subsection{Confidence intervals and multiplicity}

CIs are two-sided Student-$t$ over the $N_{\mathrm{eval}}=100$ test tasks.

\paragraph{Concern (moderate): the variance structure is mis-specified.} Each test-task mean
is itself averaged over 30 (MDP/POMDP) or 90 (APOMDP) evaluation rollouts. Treating the 100
per-task means as the only source of variance, with a $t$-interval over tasks, ignores the
within-task Monte Carlo variance and the fact that all curves are computed on the \emph{same}
100 test tasks across the $x$-axis sweep (number of training tasks). The latter induces strong
positive correlation along each curve; the apparent smoothness and the overlap/separation of
shaded bands are therefore not independent evidence point-to-point. A hierarchical / clustered
variance estimate, or at minimum paired comparisons (same test tasks, SFT vs.\ ICL vs.\ DPT),
would be the correct treatment. Paired tests would also tighten the SFT-vs-DPT comparison,
which is currently eyeballed from overlapping bars.

\paragraph{Concern (minor): no correction for the many simultaneous comparisons.} Dozens of
curves and bar groups are compared across $T$, $q$, $\alpha$, $|\mathcal{M}|$, and training-set
size with $95\%$ CIs read individually. No family-wise or FDR control is mentioned. With this
many contrasts some ``separations'' are expected by chance.

\subsection{Single-seed / single-run ambiguity}

\paragraph{Concern (moderate): fine-tuning stochasticity is not propagated.} It is not stated
whether each reported point is one QLoRA run or an average over fine-tuning seeds. Given
QLoRA's sensitivity to seed, init, and data order, and the multi-day cost acknowledged in
\S C, it is likely a single run per configuration. If so, the CIs capture
\emph{test-task} sampling variance but \emph{not} training variance---arguably the larger
source for an 7B model on a few hundred tasks. The non-monotonicities visible in several
curves (e.g.\ POMDP $q=0.5$ crossing in Fig.\ 2 right, APOMDP wiggles in Fig.\ 3) are
consistent with run-to-run noise that the CIs do not represent.

\section{Theory and its linkage to the experiments}

The theory (\S4, \S E) is a clean restatement of Zhang--Frei--Bartlett (ref.\ [8]) in
$Q$-learning notation: a single LSA layer trained by gradient flow on a population MSE
converges to the closed-form in-context estimator
$\Qhat(x_q\mid P)=x_q^\top \Gamma^{-1}\!\big(\tfrac1M\sum_i y_i x_i\big)$ with
$\Gamma=(1+\tfrac1N)\Lam+\tfrac{\mathrm{tr}\Lam}{N}I_d$, yielding a prediction-error bound (Eq.\ 6)
that the authors lift to an end-to-end policy gap (Prop.\ 1) via a finite-horizon
performance-difference identity and an on-policy calibration constant $\Con$.

\paragraph{Strength.} The decomposition into a $1/M$ in-context estimation term and a $1/N^2$
training-length bias is genuinely useful intuition, and the sample-complexity corollary
(Prop.\ 2: $\Ktest=O(\Con T d^2/\varepsilon^2)$, $K=O(\sqrt{\Con}\,d^{3/2}/\varepsilon)$)
gives a crisp, falsifiable qualitative prediction---support trajectories scale as
$1/\varepsilon^2$, training trajectories as $1/\varepsilon$. The numerical validation in \S E.2
(Fig.\ 4) of the simulation-lemma step is the right sanity check and is convincingly below the
bound across 63 configurations.

\paragraph{Concern (major): the theory is about a different object than the experiments.} The
gap between model and theory is wide enough that the theory is best read as
\emph{motivation}, and the paper should say so more plainly:
\begin{itemize}[leftmargin=1.4em,itemsep=2pt,topsep=2pt]
  \item \textbf{Architecture mismatch.} The theory is a \emph{single linear} self-attention
  layer with no softmax, no MLP, no depth, no LayerNorm, trained \emph{from scratch by
  gradient flow on population MSE}. The experiment is a 32-layer pretrained Llama-2-7B fine-tuned
  with 4-bit QLoRA on a next-token cross-entropy loss. There is no theorem connecting the QLoRA
  next-token objective to the LSA MSE minimizer, so Eq.\ (5) does not describe the trained model.
  \item \textbf{Objective mismatch.} Lemma 1 governs MSE regression onto scalar $Q$-targets.
  The model is never trained to predict $Q$-values; it is trained to predict the
  \emph{action token} of the oracle. Action imitation and $Q$-regression coincide only under
  assumptions (linear-realizable $Q$, argmax decoding, matched query covariance) that are not
  established for the tabular experiments.
  \item \textbf{Setting mismatch.} The theory requires linear-realizable $Q^*$ and a
  \emph{fixed} feature covariance $\Lam\succ0$ across train and test prompts---an assumption the
  authors correctly flag as load-bearing (covariate shift breaks the single-LSA estimator). The
  experiments are tabular, with no feature map, and POMDP/APOMDP $Q$-functions that are
  explicitly \emph{not} linear in any fixed map (\S E.4 concedes Prop.\ 1 ``does not apply
  directly''). The $\Con$ constant absorbs exactly the on-policy covariate shift the whole
  construction is fragile to, and it is never measured in the LLM experiments---only in the
  synthetic LSA simulation.
\end{itemize}
None of these are fatal to the \emph{paper}, but the abstract's ``we interpret a fine-tuned
attention layer as implicitly estimating optimal $Q$-functions'' overstates what is shown. A
single \emph{pretrained-then-fine-tuned} layer is not analyzed; a fresh single LSA layer is.
The honest claim is: ``in a stylized linear-MDP regression model, a single trained LSA layer
implements a known in-context estimator, which motivates---but does not characterize---the
behavior of the fine-tuned LLM.''

\paragraph{Concern (moderate): the headline theoretical prediction is not tested where it
would bite.} Prop.\ 2 predicts a sharp asymmetry between $\Ktest$ (number of few-shot support
trajectories) and $K$ (training rollouts per task). The experiments fix evaluation at
\emph{two} support trajectories throughout and never sweep $\Ktest$ in the main results, so
the $1/\varepsilon^2$ vs.\ $1/\varepsilon$ distinction---the most actionable theoretical
output---goes untested on the LLM. (\S F mentions an ``alternative few-shot demonstration
analysis,'' but the support-count scaling that the theory predicts is not shown in the
provided material.) Conversely the training-task sweep tests $K$ but not in a way that isolates
the $1/N^2$ bias from confounded data-diversity effects. The cleanest contribution would be a
single LLM experiment that varies $\Ktest$ and $K$ on a grid and checks the predicted slopes.

\section{Baselines}

Three references (random floor; ICL-only on the same base model; DPT [11]) plus OOD and
support-quality ablations.

\paragraph{Strength.} ICL-only on the identical base model is the correct ablation to isolate
the \emph{value of fine-tuning} from the value of the pretrained prior, and DPT is the right
prior-art comparator for in-context RL via supervised pretraining.

\paragraph{Concern (moderate): the DPT comparison is not obviously apples-to-apples on
compute or capacity.} DPT trains a transformer from scratch; the proposed method fine-tunes a
7B model carrying a massive pretrained prior. Outperforming DPT could reflect the pretrained
prior, the parameter count, or the method---these are confounded. A fairer pairing would
either (a) give DPT comparable capacity / a pretrained init, or (b) report compute-matched
results. As stated, ``our approach beats DPT'' partly restates ``a fine-tuned 7B LLM beats a
small from-scratch transformer,'' which is expected.

\paragraph{Concern (minor): missing a behavioral-cloning-without-context baseline.} A natural
and revealing baseline is the same QLoRA fine-tune evaluated with \emph{no} support
trajectories in the prompt (pure parametric policy). The delta between that and the few-shot
SFT policy is precisely the in-context contribution the paper is about, separated from the
parametric absorption of the task family. Its absence makes it hard to tell how much of the
test-time behavior is genuine in-context adaptation vs.\ a fixed policy memorized across the
one-parameter family (cf.\ \S2.1).

\paragraph{Strength.} The support-quality ablation (\S F.3: optimal vs.\ random few-shot
demonstrations) is the most diagnostic experiment in the paper. The large $T=5$ gap
($\sim$3--5\% vs.\ $\sim$20\%) is real evidence that the model \emph{uses} the context content,
not just its format---directly countering the memorization worry above for short horizons. It
would be even stronger if extended to the POMDP/APOMDP settings and tied to the $\Ktest$ sweep.

\section{Reproducibility}

\paragraph{Strength.} Hyperparameters (\S C) are reported in useful detail: QLoRA NF4 +
bf16, rank $r=64$, $\alpha=64$, dropout $0.1$ on all attention/MLP projections, Paged AdamW
32-bit, lr $2\times10^{-4}$ cosine with $0.05$ warmup, 3 epochs, A100-80GB, with an honest
wall-clock and total-project-cost disclosure. Licenses are stated. Pseudocode (\S D) is
complete enough to re-implement.

\paragraph{Concern (moderate): key choices are unstated.} (i) Number of fine-tuning seeds /
runs per point (see \S3.3 above); (ii) the exact numeric encoding / tokenizer handling of
observations and rewards; (iii) decoding temperature and how ties / invalid tokens are handled;
(iv) validation/early-stopping criterion beyond ``evaluate each epoch''; (v) whether the random
and ICL baselines use the identical prompt schema. (vi) No code or data release is mentioned in
the provided text. For a synthetic study these are all cheaply fixable and would make the
results fully reproducible.

\section{Minor and presentational points}

\begin{itemize}[leftmargin=1.4em,itemsep=3pt,topsep=2pt]
  \item \textbf{Base-model version drift.} The abstract/intro speak generically of ``pretrained
  LLMs''; \S C pins Llama-2-7B. Given the 2026 dateline this is a conservative choice and the
  results may not transfer to instruction-tuned or larger models---worth a sentence.
  \item \textbf{APOMDP ambiguity construction.} The Dirichlet-perturbation + $\mathrm{KL}\le0.2$
  entropy-ball is a reasonable but arbitrary ambiguity model; sensitivity to the $0.2$ radius is
  not shown, yet the $\alpha\times|\mathcal{M}|$ interaction is a headline result.
  \item \textbf{``Halves the gap.''} This phrase appears in the abstract and \S1 but the factor
  varies by $T$ (e.g.\ $43.2\%\to15\%$ at $T=15$ is closer to a third). Precise figures already
  in the text are stronger than the round claim.
  \item \textbf{Figure legibility.} Several multi-panel figures (Figs.\ 2--3, 5--8) pack three
  curves $\times$ three panels with overlapping CIs; the qualitative claims would be better
  supported by a companion table of point estimates $\pm$ CI at the terminal training-set size.
  \item \textbf{Darkroom (\S F.4).} Reaching $\sim$95\% of oracle on held-out goals with a
  sparse-reward $T=100$ grid is the most impressive single result and arguably deserves
  promotion from the appendix; note however the training mix (10 oracle + 100 random rollouts)
  and 5-rollout evaluation differ from the main protocol, so it is not directly comparable.
\end{itemize}

\section{Overall assessment}

\paragraph{What is solid.} The pipeline is sensible and clearly described; the ICL-only and
support-quality ablations are well chosen and genuinely diagnostic; the OOD probes add value;
reproducibility of hyperparameters is good; and the theory, taken on its own terms as a
restatement-plus-lifting of [8], is correct and the numerical check (Fig.\ 4) is appropriate.
The extension of the in-context-RL evaluation from MDPs to POMDPs/APOMDPs is a reasonable
empirical contribution.

\paragraph{What needs work, by severity.}
\begin{description}[leftmargin=2em,style=nextline,itemsep=3pt,topsep=2pt]
  \item[\textnormal{\textit{Major.}}] The theory--experiment gap (architecture, objective, and
  setting all differ) means the theoretical section motivates rather than explains the method;
  the framing should be softened accordingly, and the one cheap experiment that would partly
  bridge it---an LLM sweep over $\Ktest$ and $K$ to test the predicted scaling asymmetry---is
  absent.
  \item[\textnormal{\textit{Moderate.}}] (a) The task family is one-dimensional, so in-context
  task \emph{identification} is not cleanly separated from amortized memorization (no
  no-context fine-tuned baseline). (b) The optimality-gap metric has a fragile normalization and
  is not comparable across settings/horizons; a regret or range-normalized gap is preferable.
  (c) The variance/CI treatment ignores within-task Monte Carlo variance, curve-wise correlation,
  fine-tuning-seed variance, and multiplicity. (d) For APOMDPs ``optimal'' is the imitated oracle
  itself, so robustness is asserted but not tested under held-out ambiguity-set members. (e) The
  DPT comparison confounds method with pretrained-prior and capacity.
  \item[\textnormal{\textit{Minor.}}] Numeric serialization/tokenization unexamined; greedy
  action decoding unquantified; several unstated reproducibility details; ``halves the gap''
  imprecise.
\end{description}

\paragraph{Recommendation.} The core empirical claim---that QLoRA SFT on oracle-labeled
trajectories yields a competent few-shot in-context decision policy that beats ICL-only and
DPT on these synthetic suites---is adequately supported for short and moderate horizons and is
a useful data point. The strongest, cheapest improvements would be: (1) a richer task family
plus a no-context fine-tuned baseline to nail down that in-context adaptation (not memorization)
is doing the work; (2) a range-normalized regret metric with a hierarchical variance estimate
and at least three fine-tuning seeds per point; (3) an explicit $\Ktest$/$K$ scaling experiment
on the LLM to connect to Prop.\ 2; and (4) softened theory framing. With those, the contribution
would be both more defensible and more informative.

\end{document}
